%%%%%%%%%%%%%%%%%%%%%%%%%%%%%%%%%%%%%%%%%%%%%%%%%%%%%%%%%%%%%%%%%%%
%%% Documento LaTeX 																						%%%
%%%%%%%%%%%%%%%%%%%%%%%%%%%%%%%%%%%%%%%%%%%%%%%%%%%%%%%%%%%%%%%%%%%
% Título:	Lista de acrónimos
% Autor:  Ignacio Moreno Doblas
% Fecha:  2014-02-01, actualizado 2019-11-11
%%%%%%%%%%%%%%%%%%%%%%%%%%%%%%%%%%%%%%%%%%%%%%%%%%%%%%%%%%%%%%%%%%%
% !TEX root = A0.MiTFG.tex

%Lista de acrónimos
\chapterbeginx{Acrónimos}

\begin{acronym}[DLMS/COSEMM]
%Añade los acrónimos por orden alfabético

	\acro{ETSIT}{Escuela Técnica Superior de Ingeniería de Telecomunicación}\label{acro:ETSIT}
	\acro{TFG}{Trabajo Fin de Grado}
	\acro{UMA}{Universidad de Málaga}
        \acro{ALU}{Arithmetic Logic Unit, unidad aritmético-lógica}
        \acro{CPU}{Central Processing Unit, unidad central de procesamiento}
        \acro{CISC}{Complex Instruction Set Computing, ordenador con conjunto de instrucciones complejas}
        \acro{RISC}{Reduced Instruction Set Computing, ordenador con conjunto de instrucciones reducidas}
        \acro{FPGA}{Field Programmable Gate Array, matriz de puertas lógicas programable en campo}
        \acro{GPIO}{General Purpose Input/Output, entrada/salida de propósito general}
        \acro{I2C}{Inter-Integrated Circuit, circuito inter-integrado}
        \acro{ISA}{Instruction Set Architecture, arquitectura del conjunto instrucciones}
        \acro{LUT}{Look-Up Table, tabla de consulta}
        \acro{opcode}{Operation Code, código de operación}
        \acro{UART}{Universal Asynchronous Receiver/Transmitter, transmisor/receptor asíncrono universal}
        \acro{CPI}{Cycles Per Instruction}
        \acro{}{}
\end{acronym}

\chapterend
